\documentclass[aspectratio=169,10pt]{beamer}

% Clase 6 - Robótica Industrial
% Compilar con XeLaTeX o LuaLaTeX.
% Si se usan bloques minted, compilar con -shell-escape.

\input{~/Dropbox/Latex/TemplateSlides/template-slide.tex}

% =========================================================
% CONFIGURACIÓN DEL CURSO
% =========================================================
\renewcommand{\CourseName}{Robótica Industrial}
\renewcommand{\CourseCode}{ING 01339 / SOL-ING 220}
\renewcommand{\Faculty}{Facultad de Ingeniería}
\renewcommand{\Institution}{Politécnico Colombiano Jaime Isaza Cadavid}
\renewcommand{\Instructor}{Deimer Andrés Miranda Montoya, MSc.(c).}
\renewcommand{\Email}{deimer\_miranda91162@elpoli.edu.co}
\renewcommand{\UseImages}{1}

% =========================================================
% PIE DE PÁGINA
% =========================================================
\setbeamertemplate{footline}{%
    \leavevmode
    \hbox to \paperwidth{%
        \hspace{0.35cm}%
        \makebox[0.22\paperwidth][l]{%
            {\tiny
                \Acrobatmenu{PrevPage}{\textcolor{black}{$\triangleleft$}}%
                \hspace{0.22cm}%
                \insertpagenumber/\insertpresentationendpage%
                \hspace{0.22cm}%
                \Acrobatmenu{NextPage}{\textcolor{black}{$\triangleright$}}%
            }%
        }%
        \hfill%
        \makebox[0.70\paperwidth][r]{%
            {\tiny\Instructor}%
            \hspace{0.18cm}%
            \raisebox{-0.10cm}{\smartlogo[0.34cm]{Template/LogoPoli.png}}%
        }%
        \hspace{0.35cm}%
    }%
    \vspace{0.08cm}
}

\title{Robótica Industrial}
\subtitle{Clase 6: Composición de rotaciones en 3D}
\author{\Instructor}
\institute{\Institution}
\date{2026-2}

\begin{document}

% =========================================================
\TitleSlide

% =========================================================
\begin{frame}{Clase 6: propósito}
    \small
    \begin{columns}[T]
        \begin{column}{0.55\textwidth}
            \Idea[Propósito]{
                Comprender cómo se combinan varias rotaciones en el espacio tridimensional y por qué el orden de multiplicación cambia el resultado final.
            }

            \vspace{0.15cm}
            Al finalizar la clase, el estudiante debe poder:
            \begin{itemize}
                \item interpretar una rotación compuesta como una secuencia;
                \item reconocer qué rotación actúa primero en un producto matricial;
                \item verificar la no conmutatividad de las rotaciones 3D;
                \item distinguir rotaciones respecto a ejes fijos y ejes móviles;
                \item construir una matriz de orientación final;
                \item validar resultados mediante Python.
            \end{itemize}
        \end{column}

        \begin{column}{0.40\textwidth}
            \centering
            \begin{tikzpicture}[scale=0.95,>=Latex,every node/.style={font=\scriptsize}]
                \coordinate (O) at (0,0);
                \draw[->,very thick,BrandBlue] (O) -- (2.6,0) node[right] {$x$};
                \draw[->,very thick,BrandGreen] (O) -- (-1.0,-0.9) node[left] {$y$};
                \draw[->,very thick,BrandYellow!80!black] (O) -- (0,2.4) node[above] {$z$};
                \draw[->,thick,gray] (0.7,0) arc[start angle=0,end angle=55,radius=0.7];
                \draw[->,thick,gray] (-0.25,0.9) arc[start angle=125,end angle=45,radius=0.8];
                \node[gray] at (0.9,0.42) {$R_1$};
                \node[gray] at (0.55,1.25) {$R_2$};
                \fill (O) circle (1.4pt);
            \end{tikzpicture}

            \Result[Idea central]{
                En 3D, cambiar el orden de las rotaciones puede cambiar la orientación final.
            }
        \end{column}
    \end{columns}
\end{frame}

% =========================================================
\section{Punto de partida}
\SectionSlide{Punto de partida}{De rotaciones elementales a rotaciones compuestas}

% =========================================================
\begin{frame}{Lo que ya sabemos}
    \small
    En la clase anterior construimos las rotaciones elementales:

    \[
    R_x(\theta)=
    \begin{bmatrix}
        1&0&0\\
        0&\cos\theta&-\sin\theta\\
        0&\sin\theta&\cos\theta
    \end{bmatrix}
    \]
\vspace{0.3cm}
    \[
    R_y(\theta)=
    \begin{bmatrix}
        \cos\theta&0&\sin\theta\\
        0&1&0\\
        -\sin\theta&0&\cos\theta
    \end{bmatrix},
    \quad
    R_z(\theta)=
    \begin{bmatrix}
        \cos\theta&-\sin\theta&0\\
        \sin\theta&\cos\theta&0\\
        0&0&1
    \end{bmatrix}.
    \]

    \Result[Propiedad común]{
        Cada matriz representa una rotación rígida: conserva longitudes y ángulos, y cumple $R^TR=I$ y $\det(R)=1$.
    }
\end{frame}

% =========================================================
\begin{frame}{La nueva pregunta}
    \small
    En un robot real rara vez aparece una sola rotación aislada.

    \vspace{0.15cm}
    Una orientación puede resultar de varias acciones consecutivas:

    \[
    \text{girar alrededor de }x
    \quad\rightarrow\quad
    \text{girar alrededor de }y
    \quad\rightarrow\quad
    \text{girar alrededor de }z.
    \]

    \vspace{0.15cm}
    Esto conduce a productos matriciales como:

    \[
    R=R_z(\gamma)R_y(\beta)R_x(\alpha).
    \]

    \Warning[Pregunta clave]{
        ¿El resultado sería el mismo si escribimos $R_x(\alpha)R_y(\beta)R_z(\gamma)$?
    }

%    \Practice[Predicción]{
%        Antes de calcular, discuta si el orden de las rotaciones debería importar físicamente.
%    }
\end{frame}

% =========================================================
\begin{frame}{Composición: una acción después de otra}
    \small
    Si primero se aplica una rotación $R_1$ a un vector $\boldsymbol{p}$, se obtiene:

    \[
    \boldsymbol{p}_1=R_1\boldsymbol{p}.
    \]

    Si después se aplica una segunda rotación $R_2$:

    \[
    \boldsymbol{p}_2=R_2\boldsymbol{p}_1.
    \]

    Sustituyendo:

    \[
    \boxed{
    \boldsymbol{p}_2=R_2R_1\boldsymbol{p}}
    \]

    \begin{columns}[T]
        \begin{column}{0.48\textwidth}
            \Result[Lectura algebraica]{
                La matriz más cercana al vector actúa primero.
            }
        \end{column}
        \begin{column}{0.48\textwidth}
            \Warning[No leer de izquierda a derecha]{
                En $R_2R_1\boldsymbol{p}$, primero actúa $R_1$ y luego $R_2$.
            }
        \end{column}
    \end{columns}
\end{frame}

% =========================================================
\begin{frame}{Diagrama de composición sobre un vector}
    \small
    \begin{center}
        \resizebox{0.99\textwidth}{!}{%
    \begin{tikzpicture}[
        node distance=0.75cm,
        >=Latex,
        box/.style={draw,rounded corners,thick,minimum width=2.6cm,minimum height=0.85cm,align=center,font=\scriptsize}
    ]
        \node[box,draw=BrandBlue,fill=IdeaBody] (p) {$\boldsymbol{p}$\\vector inicial};
        \node[box,draw=BrandGreen,fill=PracticeBody,right=of p] (r1) {$R_1$\\primera rotación};
        \node[box,draw=BrandBlue,fill=IdeaBody,right=of r1] (p1) {$\boldsymbol{p}_1$\\resultado intermedio};
        \node[box,draw=BrandGreen,fill=PracticeBody,right=of p1] (r2) {$R_2$\\segunda rotación};
        \node[box,draw=BrandYellow!80!black,fill=BrandYellow!15,right=of r2] (p2) {$\boldsymbol{p}_2$\\resultado final};

        \draw[->,thick] (p) -- (r1);
        \draw[->,thick] (r1) -- (p1);
        \draw[->,thick] (p1) -- (r2);
        \draw[->,thick] (r2) -- (p2);
    \end{tikzpicture}}
    \end{center}

    \[
    \boldsymbol{p}
    \xrightarrow{R_1}
    \boldsymbol{p}_1
    \xrightarrow{R_2}
    \boldsymbol{p}_2
    \]

    \[
    \boxed{\boldsymbol{p}_2=R_2R_1\boldsymbol{p}}
    \]

    \Idea[Lectura correcta]{
        El producto matricial acumula transformaciones, pero la operación sobre el vector se evalúa desde la derecha hacia la izquierda.
    }
\end{frame}

% =========================================================
\section{El orden importa}
\SectionSlide{El orden importa}{La no conmutatividad de las rotaciones 3D}

% =========================================================
\begin{frame}{¿Qué significa no conmutatividad?}
    \small
    Para números reales:

    \[
    2\cdot3=3\cdot2.
    \]

    Pero para matrices, en general:

    \[
    AB\neq BA.
    \]

    En rotaciones 3D esto no es solamente una propiedad algebraica: tiene significado físico.

    \vspace{0.15cm}
    Dos secuencias diferentes pueden producir orientaciones finales diferentes:

    \[
    R_y(90^\circ)R_x(90^\circ)
    \neq
    R_x(90^\circ)R_y(90^\circ).
    \]

    \Warning[Idea central]{
        En 3D, girar primero alrededor de $x$ y luego alrededor de $y$ no necesariamente equivale a girar primero alrededor de $y$ y luego alrededor de $x$.
    }
\end{frame}

% =========================================================
\begin{frame}{Experimento conceptual}
    \small
    Tomemos el vector:

    \[
    \boldsymbol{p}=\begin{bmatrix}0\\0\\1\end{bmatrix}.
    \]

    Compararemos dos secuencias:

    \vspace{0.10cm}

    \begin{columns}[T]
        \begin{column}{0.48\textwidth}
            \DefBlock[Secuencia A]{
                Primero $R_x(90^\circ)$ y luego $R_y(90^\circ)$:
                \[
                \boldsymbol{p}_A=R_y(90^\circ)R_x(90^\circ)\boldsymbol{p}.
                \]
            }
        \end{column}
        \begin{column}{0.48\textwidth}
            \DefBlock[Secuencia B]{
                Primero $R_y(90^\circ)$ y luego $R_x(90^\circ)$:
                \[
                \boldsymbol{p}_B=R_x(90^\circ)R_y(90^\circ)\boldsymbol{p}.
                \]
            }
        \end{column}
    \end{columns}

    \Practice[Antes de resolver]{
        Dibuje la dirección inicial del vector $\boldsymbol{p}$ y prediga si $\boldsymbol{p}_A$ y $\boldsymbol{p}_B$ deberían coincidir.
    }
\end{frame}

% =========================================================
\begin{frame}{Secuencia A: primero $R_x$, luego $R_y$}
    \small
    \[
    \boldsymbol{p}=\begin{bmatrix}0\\0\\1\end{bmatrix}
    \]

    Primero:

    \[
    R_x(90^\circ)\boldsymbol{p}
    = \pause
    \begin{bmatrix}
        1&0&0\\
        0&0&-1\\
        0&1&0
    \end{bmatrix}
    \begin{bmatrix}0\\0\\1\end{bmatrix}
    =
    \begin{bmatrix}0\\-1\\0\end{bmatrix}.
    \]

    Luego:

    \[
    R_y(90^\circ)
    \begin{bmatrix}0\\-1\\0\end{bmatrix}
    = \pause
    \begin{bmatrix}0\\-1\\0\end{bmatrix}.
    \]
\vspace{0.2cm}
    \[
    \boxed{\boldsymbol{p}_A=\begin{bmatrix}0\\-1\\0\end{bmatrix}}
    \]
\vspace{-0.2cm}
\pause
    \Result[Interpretación]{
        Después de la primera rotación el vector queda sobre el eje $-y$. Una rotación alrededor de $y$ no lo cambia.
    }
\end{frame}

% =========================================================
\begin{frame}{Secuencia B: primero $R_y$, luego $R_x$}
    \small
    \[
    \boldsymbol{p}=\begin{bmatrix}0\\0\\1\end{bmatrix}
    \]

    Primero:

    \[
    R_y(90^\circ)\boldsymbol{p}
    =\pause
    \begin{bmatrix}
        0&0&1\\
        0&1&0\\
        -1&0&0
    \end{bmatrix}
    \begin{bmatrix}0\\0\\1\end{bmatrix}
    =
    \begin{bmatrix}1\\0\\0\end{bmatrix}.
    \]

    Luego:

    \[
    R_x(90^\circ)
    \begin{bmatrix}1\\0\\0\end{bmatrix}
    =\pause
    \begin{bmatrix}1\\0\\0\end{bmatrix}.
    \]
\vspace{0.2cm}
    \[
    \boxed{\boldsymbol{p}_B=\begin{bmatrix}1\\0\\0\end{bmatrix}}
    \]
\vspace{-0.2cm}
\pause
    \Result[Interpretación]{
        Después de la primera rotación el vector queda sobre el eje $+x$. Una rotación alrededor de $x$ no lo cambia.
    }
\end{frame}

% =========================================================
\begin{frame}{Resultado del experimento}
    \small
    Obtuvimos:

    \[
    \boldsymbol{p}_A=\begin{bmatrix}0\\-1\\0\end{bmatrix},
    \qquad
    \boldsymbol{p}_B=\begin{bmatrix}1\\0\\0\end{bmatrix}.
    \]

    Por tanto:

    \[
    \boxed{\boldsymbol{p}_A\neq\boldsymbol{p}_B}
    \]

    \begin{center}
    \begin{tikzpicture}[scale=0.65,>=Latex,every node/.style={font=\scriptsize}]
        \coordinate (O) at (0,0);
        \draw[->,thick] (O) -- (3.0,0) node[right] {$x$};
        \draw[->,thick] (O) -- (-1.2,-1.0) node[left] {$y$};
        \draw[->,thick] (O) -- (0,2.5) node[above] {$z$};
        \draw[->,very thick,BrandBlue] (O) -- (0,2.0) node[above left] {$\boldsymbol p$};
        \draw[->,very thick,BrandGreen] (O) -- (2.3,0) node[below right] {$\boldsymbol p_B$};
        \draw[->,very thick,BrandYellow!80!black] (O) -- (1.0,0.85) node[above right] {$\boldsymbol{p_A}$};
        \node[align=center] at (1.0,-1.45) {Secuencias distintas\\producen direcciones distintas};
    \end{tikzpicture}
    \end{center}
 \vspace{-0.2cm}
    \Warning[Conclusión]{
        Las rotaciones 3D no conmutan: el orden de aplicación forma parte del significado físico de la orientación final.
    }
\end{frame}

% =========================================================
\section{Cómo leer productos de rotaciones}
\SectionSlide{Cómo leer productos de rotaciones}{Primero actúa la matriz más cercana al vector}

% =========================================================
\begin{frame}{Producto de dos rotaciones}
    \small
    Si definimos:

    \[
    R=R_y(\beta)R_x(\alpha),
    \]

    entonces para un vector:

    \[
    \boldsymbol{p}'=R\boldsymbol{p}=R_y(\beta)R_x(\alpha)\boldsymbol{p}.
    \]

    La secuencia real es:
\pause
    \[
    \boldsymbol{p}
    \xrightarrow{R_x(\alpha)}
    \boldsymbol{p}_1
    \xrightarrow{R_y(\beta)}
    \boldsymbol{p}'.
    \]

    \begin{center}
    \begin{tikzpicture}[scale=0.9,>=Latex,every node/.style={font=\scriptsize}]
        \node[draw,rounded corners,thick,fill=IdeaBody,minimum width=2.4cm,minimum height=0.75cm] (p) {$\boldsymbol p$};
        \node[draw,rounded corners,thick,fill=PracticeBody,right=0.6cm of p,minimum width=2.4cm,minimum height=0.75cm] (rx) {$R_x(\alpha)$};
        \node[draw,rounded corners,thick,fill=IdeaBody,right=0.6cm of rx,minimum width=2.4cm,minimum height=0.75cm] (p1) {$\boldsymbol p_1$};
        \node[draw,rounded corners,thick,fill=PracticeBody,right=0.6cm of p1,minimum width=2.4cm,minimum height=0.75cm] (ry) {$R_y(\beta)$};
        \node[draw,rounded corners,thick,fill=BrandYellow!15,draw=BrandYellow!80!black,right=0.6cm of ry,minimum width=2.4cm,minimum height=0.75cm] (pf) {$\boldsymbol p'$};
        \draw[->,thick] (p) -- (rx);
        \draw[->,thick] (rx) -- (p1);
        \draw[->,thick] (p1) -- (ry);
        \draw[->,thick] (ry) -- (pf);
    \end{tikzpicture}
    \end{center}

    \Result[Regla operativa]{
        En una expresión como $R_3R_2R_1\boldsymbol{p}$, el vector primero recibe $R_1$, luego $R_2$ y finalmente $R_3$.
    }
\end{frame}

% =========================================================
\begin{frame}{Producto de matrices como orientación final}
    \small
    Si dos matrices son rotaciones válidas:

    \[
    R_1\in SO(3),
    \qquad
    R_2\in SO(3),
    \]

    entonces su producto también es una rotación válida:

    \[
    R=R_2R_1\in SO(3).
    \]

    \vspace{0.15cm}
    Esto significa que:

    \[
    R^TR=I,
    \qquad
    \det(R)=1.
    \]

    \Idea[Interpretación]{
        Componer rotaciones no destruye la naturaleza rígida del movimiento. El resultado sigue siendo una orientación válida.
    }

%    \Practice[Compruebe]{
%        Use las propiedades del determinante para justificar por qué $\det(R_2R_1)=1$ si $\det(R_1)=\det(R_2)=1$.
%    }
\end{frame}

% =========================================================
\begin{frame}{Demostración breve de ortogonalidad}
    \small
    Sea:

    \[
    R=R_2R_1.
    \]

    Entonces:

    \[
    R^TR
    =
    (R_2R_1)^T(R_2R_1).
    \]

    Usando:

    \[
    (AB)^T=B^TA^T,
    \]
\pause
    se obtiene:

    \[
    R^TR
    =
    R_1^TR_2^TR_2R_1.
    \]

    Como $R_2^TR_2=I$:

    \[
    R^TR=R_1^TIR_1=R_1^TR_1=I.
    \]

    \Result[Conclusión]{
        El producto de rotaciones conserva la ortogonalidad de la matriz resultante.
    }
\end{frame}

% =========================================================
\section{Ejes fijos y ejes móviles}
\SectionSlide{Ejes fijos y ejes móviles}{La misma multiplicación con interpretaciones distintas}

% =========================================================
\begin{frame}{La dificultad conceptual}
    \small
    En 3D hay que distinguir entre dos situaciones:

    \begin{columns}[T]
        \begin{column}{0.45\textwidth}
            \DefBlock[Rotaciones respecto a ejes fijos]{
                Cada rotación se interpreta alrededor de los ejes del sistema original o global.
            }
        \end{column}
        \begin{column}{0.50\textwidth}
            \DefBlock[Rotaciones respecto a ejes móviles]{
                Cada nueva rotación se interpreta alrededor de los ejes del sistema que ya fue rotado.
            }
        \end{column}
    \end{columns}

    \vspace{0.15cm}

    \Warning[Atención]{
        Esta diferencia no es un detalle de notación. Cambia la interpretación física de la secuencia de movimientos.
    }

    \Idea[Puente robótico]{
        En un manipulador, muchas rotaciones se entienden naturalmente respecto a marcos que se van moviendo con cada eslabón.
    }
\end{frame}

% =========================================================
\begin{frame}{Rotaciones respecto a ejes fijos}
    \small
    En una secuencia de rotaciones respecto a ejes fijos, cada nueva rotación se realiza alrededor de un eje del sistema base \(\{0\}\), aunque el sistema móvil ya haya cambiado de orientación

    \begin{center}
\begin{tikzpicture}[
    scale=0.95,
    >=Latex,
    every node/.style={font=\scriptsize}
    ]

    % =========================================================
    % ESTADO INICIAL
    % =========================================================
    \coordinate (O) at (0,0);

    % Sistema base {0}
    \draw[->,very thick,BrandBlue]
    (O) -- (3.0,0)
    node[right] {$x_0$};

    \draw[->,very thick,BrandGreen]
    (O) -- (-1.1,-1.0)
    node[left] {$y_0$};

    \draw[->,very thick,BrandYellow!80!black]
    (O) -- (0,2.6)
    node[above] {$z_0$};

    \node[below right] at (O) {$\{0\}$};

    % Flecha de transición
    \draw[->,very thick,gray]
    (3.3,1.2) -- (4.3,1.2);

    \node[gray,above]
    at (3.8,1.2)
    {$R_{z_0}(\theta)$};

    % =========================================================
    % DESPUÉS DE LA ROTACIÓN
    % =========================================================
    \begin{scope}[shift={(6.0,0)}]

        \coordinate (P) at (0,0);

        % -----------------------------------------------------
        % Sistema base: permanece fijo
        % -----------------------------------------------------
        \draw[->,thick,BrandBlue!45]
        (P) -- (3.0,0)
        node[right] {$x_0$};

        \draw[->,thick,BrandGreen!45]
        (P) -- (-1.1,-1.0)
        node[left] {$y_0$};

        \draw[->,thick,BrandYellow!55!black]
        (P) -- (0,2.6)
        node[above] {$z_0$};

        % -----------------------------------------------------
        % Sistema rotado {1}
        % Rotación ilustrativa de 40 grados alrededor de z0
        % -----------------------------------------------------
        \begin{scope}[rotate=40]

            \draw[->,very thick,BrandBlue]
            (P) -- (2.7,0)
            node[right] {$x_1$};

            \draw[->,very thick,BrandGreen]
            (P) -- (-1.0,-0.9)
            node[left] {$y_1$};

        \end{scope}

        % z1 coincide con z0
        \draw[->,very thick,BrandYellow!80!black]
        (P) -- (0,2.6);

        \node[right,BrandYellow!80!black]
        at (0.08,2.15)
        {$z_1=z_0$};

        % Ángulo
        \draw[->,thick,gray]
        (0.8,0)
        arc[
        start angle=0,
        end angle=40,
        radius=0.8
        ];

        \node[gray]
        at (1.02,0.32)
        {$\theta$};

        \node[below right]
        at (P)
        {$\{1\}$};

    \end{scope}

\end{tikzpicture}
    \end{center}

\begin{columns}[T]

    \begin{column}{0.48\textwidth}
        \Result[Ejes fijos]{
            Los ejes \(x_0,y_0,z_0\) pertenecen al sistema base
            y no cambian después de la rotación.
        }
    \end{column}

    \begin{column}{0.48\textwidth}
        \Result[Sistema rotado]{
            Al aplicar \(R_{z_0}(\theta)\), los ejes
            \(x_1\) y \(y_1\) cambian de orientación,
            mientras \(z_1=z_0\).
        }
    \end{column}

\end{columns}
\end{frame}

% =========================================================
\begin{frame}{Rotaciones respecto a ejes móviles}
    \small

    En una secuencia de rotaciones respecto a ejes móviles, cada nueva
    rotación se realiza alrededor de un eje del sistema que ya fue rotado.

    \begin{center}
        \begin{tikzpicture}[
            scale=0.88,
            >=Latex,
            every node/.style={font=\scriptsize}
            ]

            % =========================================================
            % SISTEMA INICIAL {0}
            % =========================================================
            \coordinate (O0) at (0,0);

            \draw[->,very thick,BrandBlue]
            (O0) -- (2.4,0)
            node[right] {$x_0$};

            \draw[->,very thick,BrandGreen]
            (O0) -- (-0.9,-0.8)
            node[left] {$y_0$};

            \draw[->,very thick,BrandYellow!80!black]
            (O0) -- (0,2.2)
            node[above] {$z_0$};

            \node[below right] at (O0) {$\{0\}$};

            % ---------------------------------------------------------
            % Primera transformación
            % ---------------------------------------------------------
            \draw[->,very thick,gray]
            (2.8,1.25) -- (3.8,1.25);

            \node[gray,above] at (3.3,1.25)
            {$R_{z_0}(\alpha)$};


            % =========================================================
            % SISTEMA {1}
            % después de rotar alrededor de z0
            % =========================================================
            \coordinate (O1) at (5.2,0);

            % x1 rotado
            \draw[->,very thick,BrandBlue]
            (O1) -- ++(2.15,1.05)
            node[right] {$x_1$};

            % y1 rotado
            \draw[->,very thick,BrandGreen]
            (O1) -- ++(0.45,-1.15)
            node[below] {$y_1$};

            % z1 NO cambia
            \draw[->,very thick,BrandYellow!80!black]
            (O1) -- ++(0,2.2)
            node[above] {$z_1=z_0$};

            \node[below right] at (O1) {$\{1\}$};

            % arco alpha alrededor de z0
            \draw[->,thick,gray]
            ($(O1)+(0.75,0)$)
            arc[start angle=0,end angle=26,radius=0.75];

            \node[gray] at ($(O1)+(0.95,0.25)$) {$\alpha$};


            % ---------------------------------------------------------
            % Segunda transformación
            % ---------------------------------------------------------
            \draw[->,very thick,gray]
            (7.9,1.25) -- (8.9,1.25);

            \node[gray,above] at (8.4,1.25)
            {$R_{x_1}(\beta)$};


            % =========================================================
            % SISTEMA {2}
            % después de rotar alrededor de x1
            % =========================================================
            \coordinate (O2) at (10.3,0);

            % ---------------------------------------------------------
            % Ejes anteriores como referencia
            % ---------------------------------------------------------
            \draw[->,thick,dashed,BrandBlue!45]
            (O2) -- ++(2.15,1.05)
            node[right] {$x_1$};

            \draw[->,thick,dashed,BrandGreen!45]
            (O2) -- ++(-0.45,-1.15)
            node[left] {$y_1$};

            \draw[->,thick,dashed,BrandYellow!50!black]
            (O2) -- ++(0,2.2)
            node[above] {$z_1$};

            % ---------------------------------------------------------
            % x2 coincide exactamente con x1
            % ---------------------------------------------------------
            \draw[->,very thick,BrandBlue]
            (O2) -- ++(2.15,1.05)
            node[above right] {$x_2=x_1$};

            % ---------------------------------------------------------
            % y2 y z2 cambian por la rotación alrededor de x1
            % Proyección esquemática
            % ---------------------------------------------------------
            \draw[->,very thick,BrandGreen]
            (O2) -- ++(0.10,-1.30)
            node[below] {$y_2$};

            \draw[->,very thick,BrandYellow!80!black]
            (O2) -- ++(-0.75,1.95)
            node[above left] {$z_2$};

            \node[below right] at (O2) {$\{2\}$};

            % indicación del giro alrededor de x1
            \draw[->,thick,gray]
            ($(O2)+(0.55,0.28)$)
            arc[start angle=25,end angle=105,
            x radius=0.52,y radius=0.30];

            \node[gray] at ($(O2)+(0.35,0.78)$) {$\beta$};

        \end{tikzpicture}
    \end{center}

    \begin{columns}[T]

        \begin{column}{0.48\textwidth}
            \Result[Primer giro]{
                Al aplicar \(R_{z_0}(\alpha)\),
                \(z_1=z_0\). Solamente cambian
                las direcciones de \(x\) y \(y\).
            }
        \end{column}

        \begin{column}{0.48\textwidth}
            \Result[Segundo giro]{
                La segunda rotación se realiza alrededor de \(x_1\).
                Por tanto, \(x_2=x_1\), mientras \(y_1\) y \(z_1\)
                cambian de dirección.
            }
        \end{column}

    \end{columns}

\end{frame}

% =========================================================
\begin{frame}{Pre-multiplicar y post-multiplicar}
    \small
    Una forma práctica de recordar la diferencia es:

    \begin{columns}[T]
        \begin{column}{0.48\textwidth}
            \DefBlock[Pre-multiplicación]{
                \[
                R_{nuevo}=R_{adicional}R_{actual}
                \]
                La rotación adicional se interpreta respecto a ejes fijos o globales.
            }
        \end{column}
        \begin{column}{0.48\textwidth}
            \DefBlock[Post-multiplicación]{
                \[
                R_{nuevo}=R_{actual}R_{adicional}
                \]
                La rotación adicional se interpreta respecto a ejes móviles o locales.
            }
        \end{column}
    \end{columns}

    \vspace{0.15cm}
    \Warning[Convención]{
        Esta regla depende de la convención de vectores columna que estamos utilizando. Por eso es indispensable mantener una notación consistente.
    }
\end{frame}

% =========================================================
\begin{frame}{Ejemplo conceptual: eje global frente a eje local}
    \small
    Suponga que un marco primero gira $90^\circ$ alrededor de $z_0$.

    \vspace{0.10cm}
    Después se solicita una rotación alrededor de $x$.

    \begin{columns}[T]
        \begin{column}{0.48\textwidth}
            \DefBlock[Si el eje es fijo]{
                La rotación se hace alrededor de $x_0$, aunque el marco ya esté rotado.
            }
        \end{column}
        \begin{column}{0.48\textwidth}
            \DefBlock[Si el eje es móvil]{
                La rotación se hace alrededor de $x_1$, que ya no coincide necesariamente con $x_0$.
            }
        \end{column}
    \end{columns}

    \vspace{0.15cm}
    \Practice[Discusión]{
        Dibuje $x_0$ y $x_1$ después de una rotación de $90^\circ$ alrededor de $z_0$. ¿Son el mismo eje?
    }
\end{frame}

% =========================================================
\section{Ejemplos resueltos}
\SectionSlide{Ejemplos resueltos}{Calcular, interpretar y comprobar}

% =========================================================
\begin{frame}{Ejemplo 1: producto de dos rotaciones}
    \small
    Construyamos:

    \[
    R=R_y(90^\circ)R_x(90^\circ).
    \]

    Sabemos que:

    \[
    R_x(90^\circ)=
    \begin{bmatrix}
        1&0&0\\
        0&0&-1\\
        0&1&0
    \end{bmatrix}
    \]
\vspace{0.2cm}
    \[
    R_y(90^\circ)=
    \begin{bmatrix}
        0&0&1\\
        0&1&0\\
        -1&0&0
    \end{bmatrix}.
    \]

    \Practice[Antes de multiplicar]{
        ¿Cuál matriz actúa primero sobre un vector si escribimos $R_yR_x\boldsymbol{p}$?
    }
\end{frame}

% =========================================================
\begin{frame}{Ejemplo 1: resultado}
    \small
    \[
    R=R_y(90^\circ)R_x(90^\circ)
    =
    \begin{bmatrix}
        0&1&0\\
        0&0&-1\\
        -1&0&0
    \end{bmatrix}.
    \]

    Leamos las columnas:

    \[
    \boldsymbol{r}_1=\begin{bmatrix}0\\0\\-1\end{bmatrix},
    \qquad
    \boldsymbol{r}_2=\begin{bmatrix}1\\0\\0\end{bmatrix},
    \qquad
    \boldsymbol{r}_3=\begin{bmatrix}0\\-1\\0\end{bmatrix}.
    \]

    \Result[Interpretación]{
        \pause El eje $x$ del marco final apunta hacia $-z$, el eje $y$ apunta hacia $+x$ y el eje $z$ apunta hacia $-y$.
    }
\end{frame}

% =========================================================
\begin{frame}{Ejemplo 2: invertir el orden}
    \small
    Ahora calculemos:

    \[
    \tilde{R}=R_x(90^\circ)R_y(90^\circ).
    \]

    El resultado es:

    \[
    \tilde{R}=
    \begin{bmatrix}
        0&0&1\\
        1&0&0\\
        0&1&0
    \end{bmatrix}.
    \]

    Comparando:

    \[
    R_y(90^\circ)R_x(90^\circ)
    =
    \begin{bmatrix}
        0&1&0\\
        0&0&-1\\
        -1&0&0
    \end{bmatrix}
    \]
\vspace{0.2cm}
    \[
    R_x(90^\circ)R_y(90^\circ)
    =
    \begin{bmatrix}
        0&0&1\\
        1&0&0\\
        0&1&0
    \end{bmatrix}.
    \]

    \Warning[Conclusión]{\pause
        Las matrices son diferentes. La orientación final también es diferente.
    }
\end{frame}

% =========================================================
\begin{frame}{Ejemplo 3: aplicar la rotación compuesta}
    \small
    Sea:

    \[
    \boldsymbol{p}=\begin{bmatrix}1\\2\\3\end{bmatrix}
    \]

    y

    \[
    R=R_y(90^\circ)R_x(90^\circ)
    =\pause
    \begin{bmatrix}
        0&1&0\\
        0&0&-1\\
        -1&0&0
    \end{bmatrix}.
    \]

    Entonces:

    \[
    \boldsymbol{p}'=R\boldsymbol{p}
    =\pause
    \begin{bmatrix}
        2\\-3\\-1
    \end{bmatrix}.
    \]

    \Practice[Verifique]{
        Compruebe que $\|\boldsymbol{p}'\|=\|\boldsymbol{p}\|=\pause\sqrt{14}$.
    }
\end{frame}

% =========================================================
\begin{frame}{Ejemplo 4: regresar usando la transpuesta}
    \small
    Si:

    \[
    \boldsymbol{p}'=R\boldsymbol{p},
    \]

    entonces podemos recuperar el vector original usando:

    \[
    \boldsymbol{p}=R^T\boldsymbol{p}'.
    \]

    Para el ejemplo anterior:

    \[
    R^T=
    \begin{bmatrix}
        0&0&-1\\
        1&0&0\\
        0&-1&0
    \end{bmatrix}
    \]
\vspace{0.2cm}
    \[
    \boldsymbol{p}=
        \begin{bmatrix}
        0&0&-1\\
        1&0&0\\
        0&-1&0
    \end{bmatrix}
    *
    \begin{bmatrix}
        2\\-3\\-1
    \end{bmatrix}
    =\pause
    \begin{bmatrix}
        1\\2\\3
    \end{bmatrix}.
    \]

    \Result[Interpretación]{
        La transpuesta deshace la orientación compuesta porque $R^{-1}=R^T$.
    }
\end{frame}

% =========================================================
\section{Problemas para clase}
\SectionSlide{Problemas para clase}{Razonar antes de usar Python}

% =========================================================
\begin{frame}{Problema 1: orden de aplicación}
    \small
    Dada la expresión:

    \[
    \boldsymbol{p}'=R_z(30^\circ)R_y(45^\circ)R_x(60^\circ)\boldsymbol{p},
    \]

    responda:

    \begin{enumerate}
        \item ¿Cuál rotación actúa primero sobre $\boldsymbol{p}$?
        \item ¿Cuál actúa de segunda?
        \item ¿Cuál actúa al final?
        \item ¿Cómo escribiría el proceso paso a paso usando $\boldsymbol{p}_1$, $\boldsymbol{p}_2$ y $\boldsymbol{p}_3$?
    \end{enumerate}

    % Solución oculta:
    % Primero Rx(60), luego Ry(45), luego Rz(30).
    % p1=Rx(60)p, p2=Ry(45)p1, p3=Rz(30)p2.
\end{frame}

% =========================================================
\begin{frame}{Problema 2: no conmutatividad}
    \small
    Calcule y compare:

    \[
    A=R_z(90^\circ)R_x(90^\circ)
    \]

    y

    \[
    B=R_x(90^\circ)R_z(90^\circ).
    \]

    Luego aplique ambas matrices al vector:

    \[
    \boldsymbol{p}=\begin{bmatrix}1\\0\\0\end{bmatrix}.
    \]

    \Practice[Pregunta]{
        ¿Los resultados representan la misma dirección final?
    }

    % Solución oculta:
    % A p = Rz Rx [1,0,0]^T = Rz [1,0,0]^T = [0,1,0]^T.
    % B p = Rx Rz [1,0,0]^T = Rx [0,1,0]^T = [0,0,1]^T.
\end{frame}

% =========================================================
\begin{frame}{Problema 3: lectura de columnas}
    \small
    Sea:

    \[
    R=
    \begin{bmatrix}
        0&0&1\\
        1&0&0\\
        0&1&0
    \end{bmatrix}.
    \]

    Sin multiplicar por ningún vector:

    \begin{enumerate}
        \item ¿hacia dónde apunta el eje $x$ del marco final?
        \item ¿hacia dónde apunta el eje $y$ del marco final?
        \item ¿hacia dónde apunta el eje $z$ del marco final?
        \item ¿la matriz puede representar una rotación propia?
    \end{enumerate}

    % Solución oculta:
    % x_final apunta a +y; y_final apunta a +z; z_final apunta a +x.
    % Columnas ortonormales y det=1, sí puede representar rotación propia.
\end{frame}

% =========================================================
\begin{frame}{Problema 4: ejes fijos o móviles}
    \small
    Un marco se encuentra inicialmente alineado con el sistema base.

    Se realizan dos acciones:

    \begin{enumerate}
        \item girar $90^\circ$ alrededor de $z$;
        \item girar $90^\circ$ alrededor del nuevo eje $x$ del marco ya rotado.
    \end{enumerate}

    \vspace{0.15cm}

    \Practice[Interpretación]{
        ¿La segunda rotación se realiza alrededor de un eje fijo o de un eje móvil?
    }

    \Result[Idea]{
        Como el segundo eje pertenece al marco ya rotado, la secuencia corresponde a una interpretación de ejes móviles.
    }
\end{frame}

% =========================================================
\section{Validación computacional}
\SectionSlide{Validación computacional}{Python como herramienta de verificación}

% =========================================================
\begin{frame}[fragile]{Python: matrices elementales}
    \small
\begin{pythoncode}
import numpy as np

def Rx(theta_deg):
    t = np.deg2rad(theta_deg)
    c = np.cos(t)
    s = np.sin(t)
    return np.array([
        [1, 0, 0],
        [0, c, -s],
        [0, s,  c] ])

def Ry(theta_deg):
    t = np.deg2rad(theta_deg)
    c = np.cos(t)
    s = np.sin(t)
    return np.array([
        [ c, 0, s],
        [ 0, 1, 0],
        [-s, 0, c] ])
\end{pythoncode}
\end{frame}

\begin{frame}[fragile]{Python: matrices elementales}
    \small
    \begin{pythoncode}

def Rz(theta_deg):
    t = np.deg2rad(theta_deg)
    c = np.cos(t)
    s = np.sin(t)
    return np.array([
        [c, -s, 0],
        [s,  c, 0],
        [0,  0, 1]
    ])
\end{pythoncode}
\end{frame}

% =========================================================
\begin{frame}[fragile]{Python: demostrar no conmutatividad}
    \small
\begin{pythoncode}
p = np.array([0.0, 0.0, 1.0])

RA = Ry(90) @ Rx(90)
RB = Rx(90) @ Ry(90)

pA = RA @ p
pB = RB @ p

print("RA =\n", RA)
print("RB =\n", RB)
print("pA =", pA)
print("pB =", pB)
print("RA == RB?", np.allclose(RA, RB))
print("pA == pB?", np.allclose(pA, pB))
\end{pythoncode}

    \Result[Lectura esperada]{
        Las matrices y los vectores resultantes no coinciden. La diferencia confirma que el orden de rotación importa.
    }
\end{frame}

% =========================================================
\begin{frame}[fragile]{Python: verificar que el producto sigue siendo rotación}
    \small
\begin{pythoncode}
R = Rz(30) @ Ry(45) @ Rx(60)

print("R.T @ R =\n", R.T @ R)
print("det(R) =", np.linalg.det(R))

p = np.array([1.0, 2.0, 3.0])
p_rot = R @ p

print("norma inicial =", np.linalg.norm(p))
print("norma final =", np.linalg.norm(p_rot))
\end{pythoncode}

    \Practice[Verificación]{
        Compruebe que $R^TR$ es aproximadamente la identidad, que $\det(R)$ es aproximadamente 1 y que la norma del vector se conserva.
    }
\end{frame}

% =========================================================
\begin{frame}[fragile]{Python: visualización 3D básica}
    \small
\begin{pythoncode}
import matplotlib.pyplot as plt

fig = plt.figure()
ax = fig.add_subplot(111, projection="3d")
p = np.array([1.0, 0.0, 0.0])
p1 = Rx(90) @ p
p2 = Ry(90) @ p1

ax.quiver(0, 0, 0, p[0],  p[1],  p[2],  label="p")
ax.quiver(0, 0, 0, p1[0], p1[1], p1[2], label="p1")
ax.quiver(0, 0, 0, p2[0], p2[1], p2[2], label="p2")

ax.set_xlim([-1, 1])
ax.set_ylim([-1, 1])
ax.set_zlim([-1, 1])
ax.set_xlabel("x")
ax.set_ylabel("y")
ax.set_zlabel("z")
ax.legend()
plt.show()
\end{pythoncode}
\end{frame}

% =========================================================
\section{Actividad de consolidación}
\SectionSlide{Actividad de consolidación}{Construir una orientación compuesta}

% =========================================================
\begin{frame}{Actividad 1: secuencia completa}
    \small
    Sea:

    \[
    \boldsymbol{p}=\begin{bmatrix}1\\1\\0\end{bmatrix}.
    \]

    Calcule manualmente:

    \[
    \boldsymbol{p}'=R_z(90^\circ)R_x(90^\circ)\boldsymbol{p}.
    \]

    Luego calcule:

    \[
    \boldsymbol{q}'=R_x(90^\circ)R_z(90^\circ)\boldsymbol{p}.
    \]

    Compare resultados.

    \Practice[Entrega]{
        Debe mostrar el cálculo paso a paso, no solamente la respuesta final.
    }
\end{frame}

% =========================================================
\begin{frame}{Actividad 2: matriz de orientación final}
    \small
    Construya la matriz:

    \[
    R=R_z(30^\circ)R_y(45^\circ)R_x(60^\circ).
    \]

    Después:

    \begin{enumerate}
        \item verifique que $R^TR=I$;
        \item verifique que $\det(R)=1$;
        \item lea sus columnas como ejes del marco final;
        \item explique qué rotación actúa primero en la expresión.
    \end{enumerate}

    \Warning[Importante]{
        No interprete la matriz final únicamente como números. Debe leerla como una orientación.
    }
\end{frame}

% =========================================================
\begin{frame}{Actividad 3: ejes móviles}
    \small
    Un marco $\{B\}$ parte alineado con $\{A\}$.

    Realice conceptualmente:

    \begin{enumerate}
        \item una rotación de $90^\circ$ alrededor de $z_B$;
        \item luego una rotación de $90^\circ$ alrededor del nuevo $x_B$.
    \end{enumerate}

    \vspace{0.15cm}

    \begin{columns}[T]
        \begin{column}{0.48\textwidth}
            \Practice[Discusión]{
                Dibuje los ejes antes y después de cada rotación.
            }
        \end{column}
        \begin{column}{0.48\textwidth}
            \Practice[Verificación]{
                Compare la orientación con la obtenida si ambos giros se hacen respecto a ejes fijos.
            }
        \end{column}
    \end{columns}
\end{frame}

% =========================================================
\section{Cierre}
\SectionSlide{Cierre}{Qué debe quedar claro antes de avanzar}

% =========================================================
\begin{frame}{Síntesis}
    \small
    \begin{itemize}
        \item Una rotación compuesta representa una secuencia de giros.
        \item En $R_2R_1\boldsymbol{p}$, primero actúa $R_1$ y luego $R_2$.
        \item Las rotaciones 3D no conmutan en general.
        \item El producto de rotaciones válidas sigue siendo una rotación válida.
        \item Una matriz compuesta también cumple $R^TR=I$, $R^{-1}=R^T$ y $\det(R)=1$.
        \item Pre-multiplicar y post-multiplicar cambia la interpretación de ejes fijos o móviles.
        \item Las columnas de la matriz final siguen representando los ejes del marco final expresados en el sistema de referencia original.
    \end{itemize}
\end{frame}

% =========================================================
\begin{frame}{Puente hacia la siguiente clase}
    \small
    Hasta ahora hemos estudiado orientación.

    \vspace{0.15cm}
    Pero un robot no sólo cambia orientación. También cambia posición.

    \vspace{0.15cm}
    Por eso, el siguiente paso será incorporar traslación:

    \[
    \text{rotación}
    \quad+
    \text{traslación}
    \quad\rightarrow\quad
    \text{transformación rígida}
    \]

    \[
    \boxed{
    {}^A\boldsymbol{p}= {}^AR_B\,{}^B\boldsymbol{p}+{}^A\boldsymbol{p}_{B}}
    \]

    \Idea[Próxima clase]{
        Este será el puente natural hacia transformaciones homogéneas y, posteriormente, cinemática directa de manipuladores.
    }
\end{frame}

\end{document}
